\chapter{Phụ lục mã nguồn chọn lọc}
\label{app:ma-nguon}

Phụ lục này trích các đoạn mã quan trọng để người đọc đối chiếu với mô tả thuật toán. Mã đầy đủ
nằm trong \repoFile{caro/}. Các đoạn dưới giữ nguyên tên biến và cấu trúc, chỉ lược bớt phần
không liên quan.

\section{Engine bitboard}
\subsection{Kiểm tra thắng}
\begin{lstlisting}[numbers=left]
def has_five(x: int) -> bool:
    for d, start in LINES:
        t = x & (x >> d)             # cặp 2
        t &= t >> (2 * d)            # dãy 4
        t &= x >> (4 * d)            # dãy 5
        if t & start:
            return True
    return False
\end{lstlisting}
Ba phép dịch và bốn phép AND cho mỗi hướng. Mặt nạ \texttt{start} loại dãy giả do lọt hàng.
Chi phí tối đa 12 phép dịch cho cả bàn, không phụ thuộc số quân.

\subsection{Đếm cửa sổ song song}
\begin{lstlisting}[numbers=left]
def window_counts(me: int, opp: int, d: int, start: int):
    free = start
    for k in range(WIN):
        free &= ~(opp >> (k * d))
    c0 = c1 = c2 = 0
    for k in range(WIN):
        p = (me >> (k * d)) & free
        carry = c0 & p
        c0 ^= p
        carry2 = c1 & carry
        c1 ^= carry
        c2 |= carry2                 # tối đa 5 = 101b, không tràn
    return c0, c1, c2, free
\end{lstlisting}
Mỗi cửa sổ có một bộ đếm 3 bit, nhưng các bit được gom theo trọng số (bit-sliced). Phép cộng có
nhớ thực hiện đồng thời cho mọi cửa sổ.

\subsection{Ô thắng và hàm lượng giá}
\begin{lstlisting}[numbers=left]
def winning_squares(me: int, opp: int) -> int:
    empty = ~(me | opp) & FULL
    res = 0
    for d, start in LINES:
        w4 = exactly(4, *window_counts(me, opp, d, start))
        for j in range(WIN):
            res |= (w4 << (j * d)) & empty
    return res

WEIGHT = (0, 1, 8, 64, 512)
def evaluate(me: int, opp: int, weight=WEIGHT) -> int:
    return side_score(me, opp, weight) - side_score(opp, me, weight)
\end{lstlisting}
\texttt{winning\_squares} là đầu vào của luật P1/P2, nước bắt buộc của Alpha-Beta và phân tích
chiến thuật. \texttt{evaluate} là hàm lượng giá của Alpha-Beta; trọng số tăng theo cấp số nhân
để ưu tiên cửa sổ nhiều quân.

\section{Alpha-Beta}
\begin{lstlisting}[numbers=left]
def _negamax(self, me, opp, key, side, depth, alpha, beta, ply, deadline):
    self._nodes += 1
    if has_five(opp):
        return -(WIN_SCORE - ply), []
    if self.use_tt:
        entry = self._tt.get(key)
        if entry is not None:
            tt_val, tt_depth, tt_flag, tt_move = entry
            if tt_depth >= depth:
                if tt_flag == EXACT:
                    return tt_val, [tt_move] if tt_move >= 0 else []
                elif tt_flag == LOWER:
                    alpha = max(alpha, tt_val)
                elif tt_flag == UPPER:
                    beta = min(beta, tt_val)
                if self.pruning and alpha >= beta:
                    return tt_val, [tt_move] if tt_move >= 0 else []
    if depth == 0 or deadline.expired():
        return self._eval_leaf(me, opp), []
    moves = list(iter_bits(bb_candidates(me, opp, self.radius)))
    if self.forced_moves:
        forced = self._forced(me, opp, moves)
        if forced:
            moves = forced
    ordered = self._order_moves(moves, me, opp, key, ply)
    orig_alpha = alpha
    best_score = -INF
    best_pv = []
    for m in ordered:
        new_opp = me | (1 << m)
        new_key = key ^ ZOBRIST[side][m] ^ SIDE_KEY
        score, pv = self._negamax(opp, new_opp, new_key, side ^ 1,
                                  depth - 1, -beta, -alpha, ply + 1, deadline)
        score = -score
        if score > best_score:
            best_score, best_pv = score, [m] + pv
        alpha = max(alpha, score)
        if self.pruning and alpha >= beta:
            self._update_killer(m, ply)
            break
    if self.use_tt and not deadline.expired():
        flag = EXACT if orig_alpha < best_score < beta else (LOWER if best_score >= beta else UPPER)
        self._tt.put(key, best_score, depth, flag, best_pv[0] if best_pv else -1)
    return best_score, best_pv
\end{lstlisting}
Đoạn này chứa đủ bốn thành phần: kiểm tra thắng, tra TT, nước bắt buộc, sắp xếp và cắt tỉa. Chú
ý thứ tự: kiểm tra thắng trước TT, vì TT không lưu trạng thái kết thúc; và không ghi TT khi hết
giờ, vì giá trị có thể chưa trọn vẹn.

\section{MCTS}
\begin{lstlisting}[numbers=left]
while True:
    if max_sims is not None and sims >= max_sims:
        break
    if deadline.expired() and sims > 0:
        break
    g = game.clone()
    node = root
    while not node.untried and node.children:
        node = max(node.children, key=lambda ch: ch.W / ch.N + c_uct * math.sqrt(math.log(node.N) / ch.N))
        g.step(node.move)
    if node.untried and not g.is_terminal():
        m = node.untried.pop(rng.randrange(len(node.untried)))
        g.step(m)
        node = node.add_child(move=m, mover=g.side ^ 1,
                              untried=candidate_moves(g, candidate_radius), key=g.key)
    z = g.result() if g.is_terminal() else rollout_fast(g.bb[0], g.bb[1], g.side, rng,
                                                        rollout_mode, rollout_radius)
    curr = node
    while curr is not None:
        curr.N += 1
        if curr.mover is not None:
            curr.W += z if curr.mover == 0 else -z
        curr = curr.parent
    sims += 1
\end{lstlisting}
Đoạn này là lõi bốn bước. \texttt{mover} lưu bên đã đi nước dẫn tới nút, để lan ngược đúng dấu.
\texttt{rollout\_fast} chỉ kiểm tra \texttt{has\_five} trên bitboard của người vừa đi, không
kiểm tra lại toàn bàn.

\section{PUCT}
\begin{lstlisting}[numbers=left]
def select_child(node, c_puct):
    sqrt_n = math.sqrt(max(1, node.N + node.vl))
    best, best_score = node.children[0], -math.inf
    for ch in node.children:
        n = ch.N + ch.vl
        q = -(ch.W + ch.vl) / n if n else 0.0
        score = q + c_puct * ch.prior * sqrt_n / (1 + n)
        if score > best_score:
            best, best_score = ch, score
    return best

@staticmethod
def _backup(path, value, release):
    for node in reversed(path):
        node.N += 1
        node.W += value
        if release:
            node.vl -= 1
        value = -value
\end{lstlisting}
\texttt{Q} của con theo góc nhìn cha là $-\frac{W+vl}{N+vl}$ vì $W$ lưu theo góc nhìn bên đi ở
con. \texttt{vl} là virtual loss, giúp các lá trong cùng lô khác nhau. \texttt{\_backup} đổi dấu
mỗi tầng.

\section{Huấn luyện}
\subsection{SL}
\begin{lstlisting}[numbers=left]
for epoch in range(epochs):
    model.train()
    for states, policies, values in train_loader:
        states, policies, values = states.to(device), policies.to(device), values.to(device)
        optimizer.zero_grad()
        pred_policies, pred_values = model(states)
        loss_p = criterion_policy(pred_policies, policies)
        loss_v = criterion_value(pred_values.squeeze(-1), values)
        loss = loss_p + c_weight * loss_v
        loss.backward()
        optimizer.step()
    if avg_val_loss < best_val_loss:
        torch.save(model.state_dict(), out_dir / "latest.pt")
\end{lstlisting}
Điểm quan trọng: \texttt{values} đã được đổi sang góc nhìn người đi trước khi vào vòng lặp
(\texttt{mover\_relative\_values}). Nếu bỏ bước này, value head học sai dấu ở thế O đi.

\subsection{REINFORCE}
\begin{lstlisting}[numbers=left]
logits, v = model(x, mask)
v = v.squeeze(1)
logp = F.log_softmax(logits / temperature, dim=1)
chosen = logp.gather(1, action[:, None]).squeeze(1)
advantage = z - v.detach()
policy_loss = -(advantage * chosen).mean()
value_loss = F.mse_loss(v, z)
entropy = -(logp.exp() * logp.masked_fill(~mask, 0.0)).sum(dim=1).mean()
loss = policy_loss + value_coef * value_loss - entropy_coef * entropy
\end{lstlisting}
\texttt{advantage} dùng \texttt{v.detach()} để baseline không chảy gradient vào policy. Entropy
được tính trên các ô hợp lệ.

\subsection{AlphaZero}
\begin{lstlisting}[numbers=left]
logits, pred = model(states, legal)
log_probs = F.log_softmax(logits, dim=1).masked_fill(~legal, 0.0)
policy_loss = -(policy * log_probs).sum(dim=1).mean()
value_loss = F.mse_loss(pred.squeeze(1), values)
loss = policy_loss + value_loss
optimizer.zero_grad(set_to_none=True)
loss.backward()
optimizer.step()
\end{lstlisting}
Khác SL ở chỗ nhãn policy là phân phối thăm $\pi$ của PUCT, không phải nhãn chuyên gia. Value
học từ kết quả ván $z$ theo góc nhìn người đi.

\section{Ghi chú đọc mã}
Ba điểm dễ nhầm khi đọc mã:
\begin{enumerate}
  \item \texttt{result()} trả theo X (+1 X thắng), còn \texttt{step()} trả reward theo người
        vừa đi. Hai quy ước này khác nhau; mọi chỗ dùng phải kiểm tra.
  \item \texttt{W} trong MCTS lưu theo góc nhìn \texttt{mover} của nút, không theo X. Khi tính
        Q cho cha, phải đổi dấu.
  \item \texttt{board\_to\_tensor} dùng góc nhìn người đang đi; kênh 2 bằng 1 nếu X đi. Nếu đổi
        quy ước này, phải đổi cả nhãn value và policy.
\end{enumerate}